



  












To understand why our sharpening objective is meaningful as a solution
concept for autoregressive language models, we observe
that when the response
$y=(y_1,\ldots,y_H)\in\cY=\cV^{H}$ is a sequence of tokens, our objective
\cref{eq:sharpening} sharpens
toward the \emph{sequence-level} arg-max response. Computing the
sequence-level arg-max for response for a language model, i.e.,
\begin{align}
  \label{eq:inference_time}
\argmax_{y_{1:H}}\piref(y_{1:H}\mid{}x)
\end{align}
is non-trivial and in general computationally intractable
\dfc{cite}. It is natural to imagine that sequence-level arg-max decoding might have
better performance on difficult reasoning tasks compared to more naive
decoding strategies such as \emph{greedy decoding} (which computes
$y_h=\argmax_{y_h\in\cV}\piref(y_h\mid{}y_{1:h-1},x)$) or beam search
\dfc{cite}. Our first result gives a stylized example where this is
indeed the case.
    \begin{example}[Benefits of sequence-level arg-max decoding]
      \label{ex:greedy}
      \dfc{Insufficiency of greedy decoding:
        \begin{itemize}
        \item First state simple/intuitive example with a couple actions/$H=2$.
        \item Then state more generic construction.
        \end{itemize}
      }
      \akcomment{Motivating example is when $\piref$ is a KL regularized optimal policy in some MDP. Then sampling from $\piref$ and greedy decoding might not be as good as MAP.}
    \end{example}

    \paragraph{Inference-time sharpening as a skyline}
Our discussion above highlights that inference-time computation
of the sequence-level arg-max response for $\piref$---or,
\emph{inference-time sharpening}---is a skyline for
the performance of any sharpening algorithm, and suggests that at
least in theory, this can lead to meaningful benefits. To validate
this empirically, we evaluate the performance of inference-time
sharpening for two stylized, yet challenging, reasoning tasks, \gametwentyfour \citep{yao2024tree} and
\stargraph \citep{bachmann2024pitfalls}.

Since exactly computing the sequence-level arg-max is intractable, we
implement inference-time sharpening by appealing to \BestofN sampling
\citet{stiennon2020learning,touvron2023llama,gao2023scaling,yang2024asymptotics}:
Given a prompt $x\in\cX$, we draw $N$ respones
$y_1,\ldots,y_n\sim\piref(\cdot\mid{}x)$, and return the response
$\yhat=\argmax_{y_i}\log\piref(y_i\mid{}x)$. The following proposition
shows that this scheme succeeds in computing an approximate arg-max
response whenever $\piref$ has sufficient coverage. To state the
result, let us define
$\Ystar(x)\ldef{}\argmax_{y\in\cY}\piref(y\mid{}x)$; we interpret
$\Ystar(x)\subset\cY$ as a set in order to accommodate non-unique
maximizers, but we will use the notation $\ystar(x)$ to indicate a
unique maximizer when it exists (i.e.,
$\Ystar(x)=\crl{\ystar(x)}$ in this case).
  \begin{proposition}[Inference-time \BestofN]
    \label{prop:bestofn_inference}
    Let a prompt $x\in\cX$ be given. For any $\deltafail\in(0,1)$ and $\gamma\in[0,1)$, as long as
    \begin{align}
      N \geq \frac{\log(\deltafail^{-1})}{\piref(\Ystar(x)\mid{}x)},
    \end{align}
    inference-time \BestofN sampling produces a response
    $
    \yhat \in \Ystar(x)$ with probability at least $1-\deltaf$.
  \end{proposition}

\paragraph{Empirical results}
  
  \cref{fig:validation} displays our findings for \gametwentyfour and
  \stargraph. We find that when $\piref$ is a pre-trained language
  model \dfc{details?}, inference-time \bestofn sampling with
  $\log\piref$ as the verifier leads to a
  meaningful increase in success probability for both tasks compared
  to greedy decoding and beam search. This suggests that sharpening is
  indeed a meaningful solution concept.


\dfc{Add footnote: If we say sharpening without qualification, we refer to training time}





    
    

    





    
            
    

    \paragraph{Training-time sharpening as amortized inference}

    Given the results above, we can view sharpening in the sense
    of our objective in \cref{eq:sharpening} as a form of
    \emph{amortized inference}. By expending extra
    computation at training time (through self-training), we aim to obtain a sharpened policy
    $\pihat$ that can compete with inference-time \BestofN sampling
    without having to expend inference-time computation.


    

\begin{remark}[Connection to maximum a-posteriori inference]
\dfc{Add references/explanation for connection to MAP/inference
  results in TCS here. Connect to diversity/diffusivity of the model
  a-la large language monkeys.}

\end{remark}


  


\section{Sharpening Algorithms for Self-Improvement}
\label{sec:algorithms}

  \dfc{
    \begin{itemize}
  \item In the remainder of the paper, we will analyze two natural
    algorithms for this task, \bestofnalg and \rlhfalg...
  \item Sketch here that asymptotically/with no errors, both algos
    converge to sharpened policy.
  \item What happens with errors/finite samples? What are the
    tradeoffs? In the remainder of
    the paper...
  \end{itemize}
}

\subsection{Self-Improvement through SFT: \bestofnalg}

\subsection{Self-Improvement through RLHF: \rlhfalg}


  
